% =========================================================================
% MODULE: INTRODUCTION & CONCEPTUAL LINEAGE
% =========================================================================

\section{Introduction}
\label{sec:introduction}

Translating high-dimensional machine learning latent spaces onto physical mediums poses a core challenge in computational aesthetics, as generative outputs often resist deterministic mechanical replication. This paper introduces \textbf{NDH-Canvas-Adapter-Enhanced}, a geometric interpretability substrate bridging these abstractions and studio practice.

\subsection{Conceptual Lineage: The Pictorial Fictions Connection}

The framework formalizes the methodology from Simon Ingram's 2026 exhibition, \textit{Pictorial Fictions – Painting from Frances Hodgkins' Latent Space} at the Dunedin Public Art Gallery~\cite{ingram2026pictorial}, which manually translated latent layers into chromatic zones. NDH-Canvas-Adapter-Enhanced automates this via:
\begin{enumerate}
    \item \textbf{From Generative Search to Bounded Physics:} Applying explicit programmatic controls to enforce geometric invariants within the Riemannian manifold.
    \item \textbf{Algorithmic Zoning Automation:} Formalizing manual color configuration via the \textit{Tessellation Bee Segmentation Core} for physical routing systems.
\end{enumerate}

\subsection{Theoretical Primitives}

The adapter establishes reliable bounds through four frameworks:
\begin{itemize}
    \item \textbf{Tessellation Bees:} Barycentric coordinate mapping into localized layout bins.
    \item \textbf{Mandala Cognition:} Radial symmetry constrained via metric tensor deformations.
    \item \textbf{Meadow Pedagogy:} Tracking sequential state evolution as parallel transport.
    \item \textbf{Manifold Goats:} Hyperbolic saddle mapping to mitigate representation crowding.
\end{itemize}
